\scanpage{275}
\mspage{225}
Il se pourrait bien que dans $\Delta^\wedge$ on ait une dualité parfaite entre ces deux espèces de notions (il faudra que j'examine la chose), peut-être $=$ dans toute catégorie $A^\wedge$ (mais dans $\Delta^\wedge$ il y a un \unsure{autom.} involutif $K\mapsto K^{\unsure{\circ|0}}$, correspondant au $X\mapsto X^\circ$ dans $\mathrm{Cat}$ \ldots). \deleted{\unsure{Après}} Et après tout, même dans le contexte des espaces topologiques les notions de propreté d'une part, d'asphéricité locale d'autre part, pour une flèche $X\unsure{-|\to}Y$, sont bien définies, et pourraient bien s'ajuster comme le font les notions correspondantes dans $\mathrm{Cat}$, et (qui sait) donner lieu à deux versions complémentaires, ou ``duales'' l'une de l'autre, d'une notion de type d'homotopie relatif.

N'en ayant pas le cœur net, je suis un peu indécis comment présenter une description axiomatique de la théorie des $\mathrm{Ho}^{\unsure{\mathcal L|L}}_W(S)$, pour une \deleted{\unsure{local}} catégorie de modèles générale $(M,W)$. Dois-je me donner d'emblée
\[
\mathrm{Liss},\mathrm{Prop}\subset\mathrm{Fl}(M),
\]
\deleted{\unsure{avec}} stables par composition, par changement de base, contenant les iso, et tels que (peut-être)
\[
\mathrm{Liss}\cap\mathrm{Prop}\subset\mathrm{Fib}_W,
\]
% archive sheet number: 275
